\scanpage{81}
\mspage{8}
un point $x\in X_K$ est entièrement déterminé par la connaissance de \deleted{$f_{x_0,x}$}
\[
f_{x_0,x}\in H^1_{\mathrm{dr}}(\Gamma,\widehat\Pi_1\underbrace{(X,x_0)}_{\widehat{\unsure{\pi|\Pi}}(x_0)}),
\]
et a fortiori par la connaissance du $1$-cocycle
\[
f_l:\Gamma\longrightarrow\widehat\Pi^{\unsure{\text{ét}}}(x_0)\qquad\text{pour un seul}
\]
$l:x_0\longrightarrow x$ --- ou ce qui revient au $=$, par la connaissance de $g_{\unsure{l^{-1}}\deleted{\unsure{x_0}}}\in H^1_g(\Gamma,\widehat\Pi_1(x_0))$, \inserted{où $l^{-1}:x\longrightarrow x_0$.}
\margin{$f_l=(g_{l^{-1}})^{-1}$}

\unsure{\P} La détermination ``explicite'' d'un $f_l$ --- \deleted{d'un $g_{l^{-1}}$} (qui caractérise donc $x$ dans le cas envisagé) est un pb très intéressant, et sûrement profond. Comme \deleted{\unsure{on}} \inserted{le} groupe $\Gamma$ bien sûr est loin d'être concret à la main, il faut préciser d'ailleurs ce qu'on peut bien pouvoir appeler une ``détermination'' (voire dét. explicite!) de $f_l$. Pour ceci, il faut sûrement considérer qu'un $u\in\Gamma$ est repéré par les invariants $f_{l_j}(u)\in\widehat\Pi_1(x_j)$ ou $g_{l_j}(u)\in\widehat\Pi_1(y_j)$, associés plus haut à un système \inserted{fini} de générateurs \inserted{$l_j:x_j\longrightarrow y_j$} du groupoïde $\Pi_1(X,I)$ ($j\in J$, \unsure{ens|un} d'indices pour les flèches génératrices, $x_j,y_j\in I$),
\margin{80}
